\documentclass[11pt]{article}
\usepackage{amsmath}
\usepackage{siunitx}
\usepackage{graphicx}
\usepackage[margin=1in]{geometry}
\usepackage{enumitem}

\newcommand{\vect}[1]{\mathbf{#1}}
\renewcommand{\thesubsection}{\thesection(\alph{subsection})}

\begin{document}

\begin{center}
\textbf{\large PHYS 2210 Electromagnetism --- Homework 6}\\
Oliver Grant \quad\textbullet\quad Due 30 October 2026
\end{center}

\section{Field of a charged ring}
A ring of radius $R$ carries uniform charge $Q$. Find the electric field on its axis at
distance $z$ from the centre.

\subsection{Setup}
By symmetry only the $z$-component survives. A charge element $dq$ contributes
\[
dE_z = \frac{1}{4\pi\varepsilon_0}\frac{dq}{z^2+R^2}\cos\theta,
\qquad \cos\theta = \frac{z}{\sqrt{z^2+R^2}}.
\]

\subsection{Result}
Integrating over the ring,
\begin{equation}
\vect{E}(z) = \frac{1}{4\pi\varepsilon_0}\frac{Qz}{(z^2+R^2)^{3/2}}\,\hat{\vect{z}}.
\end{equation}
For $z\gg R$ this reduces to the field of a point charge.

\section{Numerical example}
Take $Q = \SI{5.0}{\nano\coulomb}$, $R = \SI{10}{\centi\metre}$ and
$z = \SI{20}{\centi\metre}$. Then
\[
E = \num{8.99e9} \times \frac{\num{5.0e-9}\times 0.20}{(0.05)^{3/2}}
  \approx \SI{804}{\volt\per\metre}.
\]

\section{Magnetic field of a solenoid}
Using Amp\`ere's law with a rectangular loop,
\[
\oint \vect{B}\cdot d\vect{l} = \mu_0 I_{\text{enc}}
\quad\Longrightarrow\quad B = \mu_0 n I.
\]
\begin{enumerate}[label=(\roman*)]
  \item For $n = \SI{1200}{\per\metre}$ and $I = \SI{2.5}{\ampere}$:
        $B = \SI{3.77}{\milli\tesla}$.
  \item Doubling the current doubles the field.
  \item The field outside an ideal solenoid is zero.
\end{enumerate}

\end{document}
